% BHU PYQ -- Manually transcribed & error-corrected
\documentclass[11pt,a4paper]{article}

\usepackage[a4paper,top=2.5cm,bottom=2.5cm,left=2.8cm,right=2.8cm,headheight=14pt]{geometry}
\usepackage{amsmath,amssymb}
\usepackage[T1]{fontenc}
\usepackage[utf8]{inputenc}
\usepackage{lmodern}
\usepackage{microtype}
\usepackage{fancyhdr}
\usepackage{enumitem}
\usepackage{hyperref}
\usepackage{lastpage}
\usepackage{parskip}

\hypersetup{colorlinks=true,urlcolor=black,linkcolor=black,
  pdftitle={BPT-601: Statistical Mechanics -- BHU PYQ},pdfauthor={Banaras Hindu University}}

\pagestyle{fancy}\fancyhf{}
\renewcommand{\headrulewidth}{0.4pt}\renewcommand{\footrulewidth}{0.4pt}
\fancyhead[L]{\small   \textit{Banaras Hindu University}}
\fancyhead[R]{\small   \textit{BPT-601: Statistical Mechanics}}
\fancyfoot[C]{\small Page  \thepage\ of \pageref{LastPage}}


\newlist{parts}{enumerate}{2}
\setlist[parts,1]{label=(\alph*),leftmargin=*,itemsep=5pt,topsep=3pt}
\setlist[parts,2]{label=(\roman*),leftmargin=*,itemsep=3pt,topsep=2pt}

\newcommand{\pts}[1]{\hfill   \textit{[#1]}}


\begin{document}


\begin{center}
  {\large\bfseries BANARAS HINDU UNIVERSITY}\\[2pt]
  {
\normalsize B.Sc. (Hons.) Semester VI Examination, 2022-23}\\[6pt]
  \rule{0.8\linewidth}{0.5pt}\\[6pt]
  {\large\bfseries Paper No. BPT-601: Statistical Mechanics}\\[4pt]
  {
\normalsize Time: 3 Hours\hfill Maximum Marks: 70}
\end{center}
\rule{\linewidth}{0.4pt}
\smallskip



\noindent   \textit{Instructions: Answer all questions. Symbols have their usual meanings.}
\bigskip




\noindent   \textbf{Question 1.} Answer any   \textbf{four} of the following questions: \pts{4 $ \times$ 5 = 20}

\begin{parts}
  \item What is the ergodic hypothesis? Explain with the help of an example of a physical system.
  \item A coin is tossed 100 times. Find the probability of obtaining 50 heads. Compute the root-mean-square deviation. Explain the result. Justify any approximations used in the calculation.
  \item Estimate the room temperature specific heat of one gram of solid copper. Justify your result.
  \item Explain the meaning of the term 'Phase space'. Find the available phase space of a classical free particle, moving in a one-dimensional space between $x = -L$ and $x = +L$, and with energy between $E$ and $E + \Delta$ ($\Delta \ll E$). What are the dimensions of this phase space volume?
  \item Oxygen is a paramagnetic gas obeying Langevin's theory of paramagnetism. Its susceptibility per unit volume, at $293\, \text{K}$ and at atmospheric pressure, is $1.80  \times 10^{-4}$ MKS units. Determine its molecular magnetic moment and express it in Bohr magneton units ($\mu_B = 9.27  \times 10^{-24}\, \text{A}\cdot \text{m}^2$).
\end{parts}

\medskip\rule{\linewidth}{0.2pt}




\noindent   \textbf{Question 2.}

\begin{parts}
  \item Two random walkers start out together at the origin, each having the same probability of making a step to the left or to the right along the $x$-axis. Assuming that they make their steps simultaneously, find: (i) the probability that they meet again after $N$ steps, and (ii) the probability that they do not meet after $N$ steps. \pts{5.5}
  \item State and explain Liouville’s theorem (no proof or derivation is necessary). Verify this theorem for the density function $\rho(q, p) = C \exp(-H(q, p)/kT)$, where $C$ is a constant and $H$ is the Hamiltonian of a system of $N$ independent classical harmonic oscillators. \pts{7}
\end{parts}

\medskip\rule{\linewidth}{0.2pt}




\noindent   \textbf{Question 3.}

\begin{parts}
  \item Consider a two-level system of $N$ classical free particles in which the allowed energies of the particles are $0$ and $\epsilon$ ($\epsilon > 0$). Let $n_0$ and $n_1$ be the number of particles in the states with energies $0$ and $\epsilon$, respectively. Let the total energy of the system be $U$.
    
\begin{parts}
      \item Derive an expression for the entropy of the system.
      \item Find the temperature as a function of $U$ and make a neat sketch of the variation of $U$ as a function of $T$.
      \item Find an expression for its specific heat. Make a neat sketch of its temperature variation.
    \end{parts} \pts{12.5}
\end{parts}
\hfill   \textbf{OR}

\begin{parts}
  \item Derive the canonical partition function of an extreme relativistic gas consisting of $N$ particles. (Write down the energy-momentum relationship by neglecting the mass-energy of the particles). Derive a relationship between its pressure and its energy density. Find its entropy, and the $C_p/C_v$ ratio. \pts{12.5}
\end{parts}

\medskip\rule{\linewidth}{0.2pt}




\noindent   \textbf{Question 4.}

\begin{parts}
  \item The density of solid copper is $8.96\, \text{g/cm}^3$. The atomic weight of copper is 63.55. Assume one conduction electron per Cu atom. Explain why the conduction electrons can be treated as a degenerate Fermi gas. Find the Fermi energy (in $ \text{eV}$) of the conduction electron gas in Cu. Calculate the corresponding Fermi temperature (in Kelvin). Find the average energy per electron at zero temperature. \pts{8}
  \item Why is the electronic specific heat of conduction electrons in copper at room temperature about $1\%$ of the value expected from the equipartition law? Explain by a rough estimation. \pts{4.5}
\end{parts}
\hfill   \textbf{OR}

\begin{parts}
  \item Write down the Fermi-Dirac and Bose-Einstein distribution functions, and explain the terms in those functions. Show that they reduce to the Maxwell-Boltzmann distribution in the classical limit. Show that in the classical limit of those quantum distribution functions, the Gibbs’ paradox is automatically resolved. \pts{7}
  \item Starting from the expression of the partition function, obtain the quantum statistical distribution function for a photon gas. Explain why the chemical potential is zero in this case. \pts{5.5}
\end{parts}

\medskip\rule{\linewidth}{0.2pt}




\noindent   \textbf{Question 5.}

\begin{parts}
  \item Consider a system of two localized and independent quantum harmonic oscillators of angular frequencies $\omega$ and $3\omega$. Find the number of microstates when they share a total energy: (i) $10\,\hbar\omega$, and (ii) $12\,\hbar\omega$. \pts{6}
  \item A sample of a long-lived radionuclide gives 939 counts in 3 minutes. What is the probability that exactly 26 counts would be observed in 5 seconds? Justify any approximations that you use in the calculation. \pts{6.5}
\end{parts}

\vfill
\rule{\linewidth}{0.4pt}

\begin{center}\small
  \href{https://bhu-exam.vercel.app/}{Click here to visit our website}\\[4pt]
  \href{https://me-aryan.vercel.app/}{Click here to visit our contributor}\\[4pt]
  \href{https://quashan-me.vercel.app/}{Click here to visit our contributor}
\end{center}

\end{document}
